\chapter{Faible latence}\label{ch:iv3:low-latency}

\begin{ivtrains}
Lire du code comme le lisent la machine et les règles du langage~: ce que fait réellement un listing, pourquoi une version est plus lente que l'autre, et sur quelle garantie
s'appuie un morceau de code concurrent. La plupart des questions nomment le compilateur et le processeur dont elles dépendent, car la réponse honnête change souvent avec eux. Les
plus avancées demandent d'écrire la structure, puis de défendre chaque ligne.
\end{ivtrains}

\begin{question}[$\star$ \iqroles{QD, Dev, LLE} \iqfirm{toutes} \iqlevel{junior}]\label{iq:iv3:low-latency:1}
Compilé avec g++ pour x86-64 sous Linux, qu'affiche le \cref{lst:iv3:ll:layout}~? Quelle partie de la réponse la norme C++ fixe-t-elle, et quelle partie relève de la plateforme~?
\omcode{code/interviews-3/04-low-latency/cpp/layout.cpp}{4}{12}{Trois dispositions en mémoire.}\label{lst:iv3:ll:layout}
\end{question}

\begin{question}[$\star$ \iqroles{QD, Dev, LLE} \iqfirm{toutes} \iqlevel{junior}]\label{iq:iv3:low-latency:2}
Le \Cref{lst:iv3:ll:spinflag} est compilé avec g++ 11.4 sur x86-64, \texttt{-std=c++20 -O2}, et son auteur attend 42. Le programme est-il défini~? Pouvez-vous dire ce qu'il affiche~? Déclarer le drapeau \texttt{volatile bool} suffit-il à le corriger, et si non, qu'est-ce qui le corrige~?
\omcode{code/interviews-3/04-low-latency/cpp/spin_flag.cpp}{5}{18}{Un drapeau et une valeur.}\label{lst:iv3:ll:spinflag}
\end{question}

\begin{question}[$\star$ \iqroles{QD, Dev, LLE} \iqfirm{crypto, teneur de marché} \iqlevel{junior}]\label{iq:iv3:low-latency:3}
Avec rustc 1.97 sur x86-64, qu'affiche le \cref{lst:iv3:ll:sizes}, ligne par ligne~? Parmi les tailles des types \texttt{Option}, lesquelles Rust garantit, et lesquelles en découlent simplement par hasard~?
\omcode{code/interviews-3/04-low-latency/rust/src/bin/sizes.rs}{1}{15}{Tailles de dix types.}\label{lst:iv3:ll:sizes}
\end{question}

\begin{question}[$\star$ \iqroles{QD, Dev, LLE} \iqfirm{crypto, teneur de marché} \iqlevel{junior}]\label{iq:iv3:low-latency:4}
Le \Cref{lst:iv3:ll:spawn} se compile et affiche 300. Supprimez le mot \texttt{move} à la ligne 5. Que se passe-t-il, et à quelle étape~? Le thread est joint dès la ligne suivante~: pourquoi cela ne suffit-il pas~? Proposez une correction qui n'utilise pas de compteur de références et qui laisse \texttt{prices} utilisable dans \texttt{main} par la suite.
\omcode{code/interviews-3/04-low-latency/rust/src/bin/spawn_borrow.rs}{1}{7}{Un thread de travail.}\label{lst:iv3:ll:spawn}
\end{question}

\begin{question}[$\star$ \iqroles{LLE, Dev, QD, QT} \iqfirm{teneur de marché} \iqlevel{junior}]\label{iq:iv3:low-latency:5}
Prenez ces chiffres ronds~: un cœur à 4~GHz, 2 instructions par cycle, et un défaut de cache qui va jusqu'à la mémoire principale et coûte 90~ns. Vous disposez d'un budget de 1~microseconde entre la réception du paquet et l'envoi de l'ordre. Combien de cycles et d'instructions cela représente-t-il~? Combien de défauts de cache allant jusqu'à la mémoire principale remplissent tout le budget~? Quelle part du budget est consommée par 5 défauts, et combien d'instructions auraient pu être exécutées pendant ce temps~?
\end{question}

\begin{question}[$\star\star$ \iqroles{LLE, Dev, QD} \iqfirm{teneur de marché, hedge fund} \iqlevel{junior, confirmé}]\label{iq:iv3:low-latency:6}
Dans le \cref{lst:iv3:ll:counters}, huit threads exécutent chacun \texttt{worker} avec leur propre \texttt{t}, sur une machine dont les lignes de cache font 64~octets. Deux threads ne touchent jamais le même compteur, et chaque incrément est atomique, pourtant le passage à l'échelle est mauvais~: plus il y a de threads, plus le débit total est faible. Expliquez le mécanisme, corrigez-le, et dites ce que coûte la correction. Comment confirmeriez-vous le diagnostic~?
\omcode{code/interviews-3/04-low-latency/cpp/counters.cpp}{9}{17}{Un compteur par thread.}\label{lst:iv3:ll:counters}
\end{question}

\begin{question}[$\star\star$ \iqroles{LLE, Dev} \iqfirm{teneur de marché, hedge fund} \iqlevel{junior, confirmé}]\label{iq:iv3:low-latency:7}
Compilez le \cref{lst:iv3:ll:atomicops} avec g++ 11.4, \texttt{-std=c++20 -O2}, pour x86-64. Parmi ces huit fonctions, lesquelles ne contiennent qu'un simple déplacement ou une simple addition, et lesquelles nécessitent une instruction verrouillée ou plus coûteuse d'une autre manière~? Expliquez le schéma. Qu'est-ce qui change sur un processeur à ordre mémoire faible~?
\omcode{code/interviews-3/04-low-latency/cpp/atomic_ops.cpp}{3}{13}{Huit petites fonctions.}\label{lst:iv3:ll:atomicops}
\end{question}

\begin{question}[$\star\star$ \iqroles{LLE, Dev, QD} \iqfirm{toutes} \iqlevel{junior, confirmé}]\label{iq:iv3:low-latency:8}
Dans le \cref{lst:iv3:ll:lists}, \texttt{values} part de $0, 1, \dots, n-1$, avec $n$ valant deux millions, et les deux listes finissent par contenir les mêmes valeurs dans le même ordre. La même fonction, qui parcourt une liste et additionne ses valeurs, met beaucoup plus de temps sur \texttt{b} que sur \texttt{a}. Pourquoi~? Comment confirmeriez-vous votre explication sans aucun compteur matériel, et que modifieriez-vous~?
\omcode{code/interviews-3/04-low-latency/cpp/list_order.cpp}{41}{45}{Deux listes égales.}\label{lst:iv3:ll:lists}
\end{question}

\begin{question}[$\star\star$ \iqroles{QD, Dev, LLE} \iqfirm{toutes} \iqlevel{junior, confirmé}]\label{iq:iv3:low-latency:9}
Le \Cref{lst:iv3:ll:countnew} remplace le \texttt{operator new} global par un opérateur qui compte ses appels. Avec g++ 11.4 et sa bibliothèque standard, \texttt{-std=c++20 -O2}, x86-64, quel nombre chaque ligne affiche-t-elle~? Lesquels de ces nombres une autre bibliothèque standard pourrait-elle modifier~?
\omcode{code/interviews-3/04-low-latency/cpp/count_new.cpp}{10}{47}{Comptage des allocations.}\label{lst:iv3:ll:countnew}
\end{question}

\begin{question}[$\star\star$ \iqroles{LLE, Dev, QD} \iqfirm{teneur de marché, crypto} \iqlevel{junior, confirmé}]\label{iq:iv3:low-latency:10}
Un paquet se trouve dans un tampon d'octets, avec un en-tête qui commence 3 octets après le début. Avec g++ 11.4 sur x86-64, \texttt{-std=c++20 -O2}, le \cref{lst:iv3:ll:parse} affiche deux fois la bonne valeur. Les deux fonctions sont-elles correctes en C++20~? À quoi vous attendez-vous pour le code machine de chacune~?
\omcode{code/interviews-3/04-low-latency/cpp/parse_header.cpp}{6}{27}{Deux façons de lire un horodatage.}\label{lst:iv3:ll:parse}
\end{question}

\begin{question}[$\star\star$ \iqroles{LLE, Dev, QD, QR} \iqfirm{toutes} \iqlevel{junior, confirmé}]\label{iq:iv3:low-latency:11}
Le \Cref{lst:iv3:ll:benchbad}, compilé avec g++ 11.4, \texttt{-O2}, sur x86-64, affiche environ 70~nanosecondes pour un milliard d'itérations. Pourquoi~? Corrigez-le, et listez ce qui cloche aussi dans ce benchmark pour mesurer la latence de \texttt{work}. Si vous lisez ensuite directement le compteur d'horodatage (TSC) avec \texttt{rdtsc}, au lieu d'une horloge fournie par une bibliothèque, quelle nouvelle erreur pouvez-vous commettre~?
\omcode{code/interviews-3/04-low-latency/cpp/bench_bad.cpp}{5}{16}{Un benchmark.}\label{lst:iv3:ll:benchbad}
\end{question}

\begin{question}[$\star\star$ \iqroles{LLE, Dev} \iqfirm{teneur de marché, hedge fund} \iqlevel{junior, confirmé}]\label{iq:iv3:low-latency:12}
Avec g++ 11.4, \texttt{-std=c++20 -O2}, sur x86-64, le \cref{lst:iv3:ll:publish} affiche 42 à chaque exécution que vous tentez. Le programme est-il correct selon le modèle mémoire de C++20~? Quelque chose changerait-il sur ARM ou POWER~? Corrigez-le avec les ordres mémoire les plus faibles qui fonctionnent.
\omcode{code/interviews-3/04-low-latency/cpp/publish.cpp}{5}{20}{Publication d'une valeur.}\label{lst:iv3:ll:publish}
\end{question}

\begin{question}[$\star\star$ \iqroles{LLE, Dev} \iqfirm{teneur de marché, hedge fund} \iqlevel{junior, confirmé}]\label{iq:iv3:low-latency:13}
Avec g++ 11.4, \texttt{-O2}, sur x86-64, quatre threads verrouillent chacun ce spinlock et incrémentent 20\,000 fois un compteur \texttt{long} ordinaire, et le total est juste. Le \Cref{lst:iv3:ll:spinlock} présente encore un défaut de correction et un défaut de performance sous contention. Trouvez les deux et corrigez-les.
\omcode{code/interviews-3/04-low-latency/cpp/spinlock.cpp}{6}{15}{Un spinlock.}\label{lst:iv3:ll:spinlock}
\end{question}

\begin{question}[$\star\star$ \iqroles{LLE, Dev, QD} \iqfirm{teneur de marché, crypto} \iqlevel{confirmé}]\label{iq:iv3:low-latency:14}
Vous voulez qu'un thread publie un meilleur cours acheteur et un meilleur cours vendeur que d'autres lisent de façon atomique. Le \Cref{lst:iv3:ll:quote} est compilé avec g++ 11.4, \texttt{-std=c++20 -O2 -latomic}, sur un processeur x86-64 qui prend en charge le compare-and-swap sur 16~octets. Il signale que le type de 8~octets est sans verrou (lock-free), et que ceux de 16 et de 24~octets ne le sont pas. Que signifie cela pour un chemin critique (hot path)~? Que se passe-t-il si vous omettez \texttt{-latomic}~? Que feriez-vous à la place~?
\omcode{code/interviews-3/04-low-latency/cpp/quote_atomic.cpp}{4}{20}{Trois tailles d'atomic.}\label{lst:iv3:ll:quote}
\end{question}

\begin{question}[$\star\star$ \iqroles{QD, Dev, LLE} \iqfirm{crypto, teneur de marché} \iqlevel{junior, confirmé}]\label{iq:iv3:low-latency:15}
En Rust, un type qui est \texttt{Send} peut être déplacé vers un autre thread, et un type qui est \texttt{Sync} peut être partagé par référence entre threads. Pour chacun de ces types, dites s'il est \texttt{Send} et s'il est \texttt{Sync}, en précisant la raison~: \texttt{Rc<u64>}, \texttt{Arc<u64>}, \texttt{Arc<Cell<u64>{}>}, \texttt{Cell<u64>}, \texttt{RefCell<u64>}, \texttt{Mutex<Cell<u64>{}>}, \texttt{Mutex<Rc<u64>{}>}, \texttt{\&\textquotesingle{}static Cell<u64>}, \texttt{*mut u64}, \texttt{MutexGuard<\textquotesingle{}static, u64>}, \texttt{AtomicU64}.
\end{question}

\begin{question}[$\star\star$ \iqroles{QD, Dev, LLE} \iqfirm{crypto, teneur de marché} \iqlevel{junior, confirmé}]\label{iq:iv3:low-latency:16}
Le \Cref{lst:iv3:ll:lost} ne contient aucun \texttt{unsafe} et se compile sans aucun avertissement avec rustc 1.97. Que peut-il afficher~? Y a-t-il une course aux données (data race), ou un comportement indéfini~? Quelle est la plus petite valeur finale que les règles autorisent, et comment corrigeriez-vous le programme~?
\omcode{code/interviews-3/04-low-latency/rust/src/bin/lost_update.rs}{1}{18}{Un compteur partagé, en code sûr.\protect\hfil\protect\break}\label{lst:iv3:ll:lost}
\end{question}

\begin{question}[$\star\star$ \iqroles{QD, Dev, LLE} \iqfirm{crypto, teneur de marché} \iqlevel{junior, confirmé}]\label{iq:iv3:low-latency:17}
Un allocateur global compte chaque appel à allocate ou à reallocate, et \texttt{show} dans le \cref{lst:iv3:ll:countalloc} affiche le nombre d'appels effectués par sa fermeture. Avec rustc 1.97 en mode release, et sachant qu'un \texttt{Vec} vide qui grandit démarre à une capacité de 4 et double à chaque fois, quel nombre chaque ligne affiche-t-elle~? Parmi ces nombres, lesquels la documentation garantit-elle~?
\omcode{code/interviews-3/04-low-latency/rust/src/bin/count_alloc.rs}{25}{53}{Comptage des appels à l'allocateur.}\label{lst:iv3:ll:countalloc}
\end{question}

\begin{question}[$\star\star$ \iqroles{QD, Dev, LLE} \iqfirm{crypto, teneur de marché} \iqlevel{confirmé}]\label{iq:iv3:low-latency:18}
Compilez le \cref{lst:iv3:ll:bounds} avec rustc 1.97, \texttt{-C opt-level=3}, sur x86-64. Un indice pouvant être hors bornes nécessite normalement une comparaison et un branchement vers un panic. Parmi ces sept fonctions, lesquelles conservent encore une vérification de bornes sur l'indexation, lesquelles effectuent une autre vérification une seule fois au départ, et lesquelles n'en ont aucune~? Tout cela est-il garanti~?
\omcode{code/interviews-3/04-low-latency/rust/src/bounds.rs}{1}{36}{Sept façons d'indexer.}\label{lst:iv3:ll:bounds}
\end{question}

\begin{question}[$\star\star$ \iqroles{LLE, Dev, QD} \iqfirm{teneur de marché, hedge fund} \iqlevel{junior, confirmé}]\label{iq:iv3:low-latency:19}
Un processus lit des entrées aléatoires d'une table de 8~GiB. Supposez que le cache de traduction (le TLB) contienne 2048 entrées, qu'un défaut dans celui-ci coûte une marche dans la table des pages de 30~ns, et que cette marche coûte la même chose quelle que soit la taille de page. Pour une taille de page de 4~KiB, puis de 2~MiB, estimez quelle part de la table le TLB couvre, le taux de défaut et le coût moyen de traduction par lecture. Quelle quantité de mémoire occupent les entrées de table des pages de dernier niveau à chaque taille de page, à raison de 8~octets par entrée~? Quels sont les pièges des grandes pages (huge pages)~?
\end{question}

\begin{question}[$\star\star\star$ \iqroles{LLE, Dev, QD} \iqfirm{teneur de marché, hedge fund} \iqlevel{confirmé}]\label{iq:iv3:low-latency:20}
Deux variables \texttt{std::atomic<int>}, \texttt{x} et \texttt{y}, sont initialisées à 0. Le thread 1 exécute \texttt{x.store(1); r1 = y.load();}. Le thread 2 exécute \texttt{y.store(1); r2 = x.load();}. Les deux registres peuvent-ils finir à 0~? Ce serait possible si (a) toutes les opérations sont relâchées (relaxed), (b) les écritures sont en release et les lectures en acquire, (c) toutes les opérations ont une cohérence séquentielle (sequentially consistent)~? Pour chaque cas, répondez pour le langage C++20 et pour le matériel x86-64. Quelle instruction fait la différence en (c), et comment montreriez-vous le résultat sur une machine réelle~?
\end{question}

\begin{question}[$\star\star\star$ \iqroles{LLE, Dev} \iqfirm{teneur de marché, hedge fund, crypto} \iqlevel{junior, confirmé}]\label{iq:iv3:low-latency:21}
Au tableau blanc, en C++20, écrivez un tampon circulaire (ring buffer) de capacité fixe, pour exactement un thread producteur et un thread consommateur, sans aucun verrou. \texttt{push} renvoie false lorsque le tampon est plein, et \texttt{pop} renvoie false lorsqu'il est vide, et la capacité est une puissance de deux. Défendez ensuite l'ordre mémoire de chaque opération atomique~: qu'est-ce qui peut mal tourner si vous l'affaiblissez~?
\end{question}

\begin{question}[$\star\star\star$ \iqroles{LLE, Dev} \iqfirm{teneur de marché, hedge fund} \iqlevel{confirmé}]\label{iq:iv3:low-latency:22}
Le \Cref{lst:iv3:ll:adjacent} est une file à un producteur et un consommateur (SPSC) correcte, les deux threads étant sur des cœurs différents. Mettre \texttt{alignas(64)} sur \texttt{head}, \texttt{tail} et \texttt{slots} la rend nettement plus rapide, et une modification supplémentaire, en gardant tous les ordres mémoire inchangés, la rend encore plus rapide. Expliquez les deux mécanismes, nommez cette modification supplémentaire, et dites ce qu'elle coûte. Comment confirmeriez-vous chacun des effets~?
\omcode{code/interviews-3/04-low-latency/cpp/spsc_layout.cpp}{11}{29}{Une file avec ses champs côte à côte.}\label{lst:iv3:ll:adjacent}
\end{question}

\begin{question}[$\star\star\star$ \iqroles{LLE, Dev, QD} \iqfirm{teneur de marché, crypto} \iqlevel{confirmé}]\label{iq:iv3:low-latency:23}
Le \Cref{lst:iv3:ll:wrap} est présenté en mono-thread pour rester court. Ses indices progressent librement et reposent sur le débordement non signé, et tous les tests passent. À un million de messages par seconde, il peut commencer à perdre des données au bout d'un moment. Au bout de combien de temps, dans quelle condition exacte, et quelle est la correction la plus simple~? Est-ce \texttt{tail - head} qui pose problème~?
\omcode{code/interviews-3/04-low-latency/cpp/wrap.cpp}{4}{21}{Un tampon circulaire aux indices libres.}\label{lst:iv3:ll:wrap}
\end{question}

\begin{question}[$\star\star\star$ \iqroles{LLE, Dev, QD, QR} \iqfirm{teneur de marché, hedge fund} \iqlevel{confirmé}]\label{iq:iv3:low-latency:24}
Un thread du chemin critique doit enregistrer chaque latence, un nombre entier de nanosecondes de 1~ns à environ 1~s, de sorte que n'importe quel centile puisse être relu plus tard à 2~\% près. Il ne doit pas allouer, ne doit pas utiliser de virgule flottante, et doit prendre un temps constant par enregistrement. Concevez l'histogramme~: combien de compartiments, quelle quantité de mémoire, et quelle est la fonction d'indexation en C++20~? Pourquoi ne pas prendre un compartiment par nanoseconde~?
\end{question}

\begin{question}[$\star\star\star$ \iqroles{LLE, Dev} \iqfirm{teneur de marché, hedge fund} \iqlevel{confirmé}]\label{iq:iv3:low-latency:25}
Le \Cref{lst:iv3:ll:firsttouch} s'exécute sur x86-64 Linux avec des pages de 4~KiB, glibc, et sans activation forcée des grandes pages transparentes. Environ quels cinq nombres affiche-t-il~? Servez-vous de votre réponse pour expliquer pourquoi les premiers messages après le démarrage d'un processus peuvent être lents, et ce que vous feriez au démarrage. Sur une machine à deux sockets, où se retrouve la mémoire du vecteur si \texttt{main} le construit et qu'un thread épinglé sur l'autre socket l'utilise ensuite~?
\omcode{code/interviews-3/04-low-latency/cpp/first_touch.cpp}{16}{35}{Défauts de page mineurs, comptés.}\label{lst:iv3:ll:firsttouch}
\end{question}

\begin{question}[$\star\star\star$ \iqroles{LLE, Dev, QD, QR, QT} \iqfirm{teneur de marché, hedge fund} \iqlevel{junior, confirmé}]\label{iq:iv3:low-latency:26}
Un outil de charge envoie une requête toutes les 2~ms, mais attend la réponse à chaque requête avant d'envoyer la suivante, et mesure le temps écoulé entre l'envoi et la réponse. Le serveur est une file FIFO unique qui met 1~ms par requête. Pendant un essai de 100~secondes, il se fige une fois, pendant exactement 1~seconde. Estimez à la main ce que l'outil affiche pour le 99\up{e}~centile, pour le 99.9\up{e}~centile et pour son maximum. Estimez ensuite les mêmes centiles pour un client qui respecte le calendrier de 2~ms quoi qu'il arrive, en mesurant à partir de l'instant d'envoi prévu de chaque requête. Lequel croyez-vous, et comment modifieriez-vous l'outil~?
\end{question}

\begin{question}[$\star\star\star$ \iqroles{LLE, Dev} \iqfirm{teneur de marché, hedge fund} \iqlevel{confirmé}]\label{iq:iv3:low-latency:27}
Un port Ethernet de 10~Gb/s reçoit des données de marché en rafale, au débit de ligne (line rate), dans les plus petites trames~: 64~octets, plus 20~octets de préambule et d'intervalle inter-trames sur le câble. Combien de temps avez-vous par paquet, et combien de cycles à 3~GHz~? Supposez qu'un appel système de réception coûte 500~ns et que la pile réseau du noyau ajoute 150~ns par paquet. Un appel par paquet suffit-il~? Un appel pour 32 paquets~? Choisissez maintenant entre des sockets groupés avec attente active (busy polling), un socket \texttt{AF\_XDP} et DPDK, alors que le même port doit aussi continuer à transporter le trafic de gestion ordinaire et que l'équipe veut garder les outils standard du noyau. À quoi renonce chacune de ces options~?
\end{question}

\begin{question}[$\star\star\star$ \iqroles{LLE, Dev} \iqfirm{teneur de marché, hedge fund} \iqlevel{confirmé}]\label{iq:iv3:low-latency:28}
Sur Linux, un thread épinglé sur un cœur retiré de l'ordonnanceur avec \texttt{isolcpus} tourne à vide sur une file et n'effectue aucun appel système. La plupart des itérations durent 50~ns, mais toutes les quelques secondes l'une d'elles dure des dizaines de microsecondes, et de temps en temps l'une dépasse la milliseconde. Qu'est-ce qui peut encore interrompre ce cœur~? Donnez l'ordre dans lequel vous procéderiez aux vérifications, l'outil ou le fichier de chaque vérification, et le correctif.
\end{question}

\begin{question}[$\star\star\star$ \iqroles{QD, Dev, LLE} \iqfirm{crypto, teneur de marché} \iqlevel{confirmé}]\label{iq:iv3:low-latency:29}
Avec rustc 1.97 sur x86-64, que montrent les deux premières lignes de sortie du \cref{lst:iv3:ll:header}, et lesquelles de ces quatre tailles sont garanties~? \texttt{ts\_cast} est ensuite appelée sur une tranche commençant à 3 octets dans un tampon aligné sur 8 octets. Que se passe-t-il en mode débogage et en mode release, et que dit le langage~? Pour une valeur \texttt{p} de type \texttt{Packed}, pourquoi écrire \texttt{\{ p.ts \}} compile-t-il, alors que prendre \texttt{\&p.ts} ne compile pas~? Comment faut-il écrire \texttt{ts\_cast}~?
\omcode{code/interviews-3/04-low-latency/rust/src/bin/header.rs}{4}{31}{Dispositions en mémoire et une conversion de type.}\label{lst:iv3:ll:header}
\end{question}

\begin{question}[$\star\star\star$ \iqroles{QD, Dev, LLE} \iqfirm{crypto, teneur de marché} \iqlevel{confirmé}]\label{iq:iv3:low-latency:30}
Le \Cref{lst:iv3:ll:mailbox} est une boîte aux lettres d'un seul élément que son auteur décrit comme sûre à utiliser depuis du code sûr. Trouvez chaque manière dont du code ne contenant aucun \texttt{unsafe} peut utiliser cette API pour provoquer un comportement indéfini, ainsi que toute manière dont il peut se comporter mal sans comportement indéfini. Corrigez-le, en modifiant l'API là où c'est nécessaire.
\omcode{code/interviews-3/04-low-latency/rust/src/bin/mailbox.rs}{7}{34}{Une boîte aux lettres d'un seul élément.}\label{lst:iv3:ll:mailbox}
\end{question}

\begin{question}[$\star\star\star\star$ \iqroles{LLE, Dev} \iqfirm{teneur de marché, hedge fund} \iqlevel{confirmé}]\label{iq:iv3:low-latency:31}
Le \Cref{lst:iv3:ll:seqlock} est un schéma lecteurs-écrivain pour une cotation~: un écrivain, de nombreux lecteurs, et un lecteur recommence s'il a vu le numéro de séquence impair ou changé pendant sa lecture. Est-il correct selon le modèle mémoire du C++20~? Nommez chaque manière dont il peut échouer, que l'architecture x86-64 la masque ou non. Réécrivez les deux côtés, et justifiez chaque ordre que vous utilisez. Que coûte le schéma à l'écrivain et aux lecteurs, par rapport à un mutex~?
\omcode{code/interviews-3/04-low-latency/cpp/seqlock_naive.cpp}{9}{29}{Un numéro de séquence et une cotation.}\label{lst:iv3:ll:seqlock}
\end{question}

\begin{question}[$\star\star\star\star$ \iqroles{LLE, Dev} \iqfirm{teneur de marché, hedge fund} \iqlevel{confirmé}]\label{iq:iv3:low-latency:32}
De nombreux threads appellent \texttt{push} et \texttt{pop} sur la pile implantée dans le \cref{lst:iv3:ll:treiber}. Elle présente deux défauts distincts, et chacun peut corrompre la mémoire ou faire planter le programme sous g++ sur x86-64, quels que soient les ordres mémoire. Donnez, pour chacun, l'ordre exact des événements entre deux threads qui le montre. Comparez ensuite les façons de le corriger~: que coûte chacune~?
\omcode{code/interviews-3/04-low-latency/cpp/treiber.cpp}{6}{31}{Une pile construite sur compare-and-swap.}\label{lst:iv3:ll:treiber}
\end{question}

\begin{question}[$\star\star\star\star$ \iqroles{LLE, Dev} \iqfirm{teneur de marché, hedge fund} \iqlevel{confirmé}]\label{iq:iv3:low-latency:33}
Dites d'un algorithme qu'il est sans verrou (lock-free) si, quels que soient les threads suspendus, un thread en cours d'exécution termine toujours son opération en un nombre fini de ses propres étapes. Le \Cref{lst:iv3:ll:mpmc} présente la partie push d'une file bornée pour de nombreux producteurs et de nombreux consommateurs. Elle n'utilise aucun mutex, et \texttt{try\_pop}, non montré, en est l'image symétrique~: il ne prend une position que lorsque le champ \texttt{seq} de cette case indique que l'élément est prêt. Supposez qu'un producteur soit suspendu juste après la réussite de son compare-and-swap de la ligne 28, sur une file de 8 cases. Que peuvent encore faire les autres threads, et que voient-ils~? La file est-elle sans verrou au sens ci-dessus~? Quelle est la latence au pire cas qu'elle peut imposer à un consommateur, et comment reconcevriez-vous la file si cela est inacceptable~?
\omcode{code/interviews-3/04-low-latency/cpp/mpmc.cpp}{7}{40}{La partie push d'une file bornée.}\label{lst:iv3:ll:mpmc}
\end{question}

\begin{question}[$\star\star\star\star$ \iqroles{QD, Dev, LLE} \iqfirm{crypto, teneur de marché} \iqlevel{confirmé}]\label{iq:iv3:low-latency:34}
En Rust, sans crate externe, concevez une file bornée à un producteur et un consommateur (SPSC) dont l'API fait échouer à la compilation ces trois erreurs~: deux threads qui poussent en même temps, un pop répété alors qu'on détient encore une référence sur un élément consulté, et l'envoi, à travers la file, d'un type d'élément qui n'est pas thread-safe. Donnez les types et les signatures, écrivez \texttt{push} et un \texttt{peek} sans copie, et indiquez quelles lignes \texttt{unsafe impl} vous sont nécessaires et pourquoi. Quelles erreurs le compilateur ne vous signalera-t-il toujours pas~?
\end{question}

\begin{question}[$\star\star\star\star$ \iqroles{QD, Dev, LLE} \iqfirm{crypto, teneur de marché} \iqlevel{confirmé}]\label{iq:iv3:low-latency:35}
Le \Cref{lst:iv3:ll:myarc} est un pointeur à comptage de références atomique écrit à la main, dans lequel chaque opération sur le compteur est \texttt{Relaxed}. Quelles opérations sont correctes en l'état, et lesquelles ne le sont pas~? Donnez un ordre d'événements qui tourne mal, et corrigez-le avec les ordres les plus faibles qui fonctionnent. Pourquoi les deux bornes \texttt{T: Send + Sync} doivent-elles être présentes, et que coûte votre correction sur x86-64 et sur un processeur à ordre mémoire faible~?
\omcode{code/interviews-3/04-low-latency/rust/src/my_arc_bad.rs}{1}{37}{Un pointeur à comptage de références.}\label{lst:iv3:ll:myarc}
\end{question}
